\documentclass[10pt,a4paper]{amsart}
\usepackage[margin=19mm]{geometry}
\usepackage{amsmath,amssymb,mathtools}
\usepackage[scr=rsfs,cal=euler]{mathalfa}
\usepackage{xfrac}
\usepackage[unicode,hidelinks,bookmarks=false]{hyperref}
\newcommand{\ie}{i.e.\ }
\newcommand{\cf}{cf.\ }
\newcommand{\resp}{respectively\ }
\newcommand{\Subsection}{\subsection}
\providecommand{\nobreakdash}{\nobreak\hbox{-}}
\setlength{\emergencystretch}{2em}
\multlinegap0pt
\usepackage{mathtools}
\usepackage[shortlabels]{enumitem}
\usepackage{amssymb}
\usepackage{xcolor}
\usepackage{algorithm}\usepackage{algpseudocode}
\usepackage{float}
\newtheorem{theorem}{Theorem}
\newtheorem{assumption}{Assumption}
\newtheorem{lemma}{Lemma}
\renewcommand{\P}{\mathbb{P}}
\title{In-Context Learning for Data-Driven Censored Inventory Control}
\author{Sohom Mukherjee, Anh-Duy Pham, Richard Pibernik, Yunbei Xu}
\date{Source: arXiv:2605.14840v1 (CC BY 4.0)}
\begin{document}
\maketitle

\noindent\emph{Source and licence.} In-Context Learning for Data-Driven Censored Inventory Control. Version arXiv:2605.14840v1 (CC BY 4.0), distributed by arXiv under the Creative Commons Attribution 4.0 International licence, http://creativecommons.org/licenses/by/4.0/, as recorded in the arXiv licence metadata for this exact version and shown on the abstract page https://arxiv.org/abs/2605.14840v1. Full licence notice: \texttt{licenses/arxiv-2605.14840-CC-BY-4.0.txt}.\par\smallskip

\noindent\textit{Editorial scope.} Counted unit: the article's theorem thm:completion\_to\_ovserved (source0.tex lines 1233--1252). The article's complete original proof of it is delivered with it (source0.tex lines 3078--3118, one \texttt{proof} environment of 41 lines, which the article places away from the statement). No mathematics is written, repaired or re-derived: the statement and the proof are the article's own lines, in the article's own notation, and no retained line is substituted or re-flowed.  Retained uncounted as the source-local context the proof needs: lines 2853--2858; lines 2870--2877; lines 2884--2894; lines 2972--2982; lines 3019--3021; lines 3023--3029; lines 3044--3051.  Nothing was omitted from any retained span, so the package carries no omission record.  No retained line contains a citation, so the wrapper carries no bibliography and no bibliography entry.  No retained line contains a figure environment, an image-inclusion command or drawing code, so no image is delivered and the package declares no asset dependency.  Editorial formatting only, in addition: an AMS \texttt{amsart} preamble that restates the article's own declarations and macros used by the retained lines, namely the environments \texttt{theorem}, \texttt{assumption}, \texttt{lemma} on separate counters, together with the standard packages those lines need.  The retained environments print in this document as Theorem 1, Assumption 1, Lemma 1 in order of appearance, whereas the article prints the same items with the numbering of its own full document.  The complete native source of the article as distributed in the arXiv e-print for this exact version is bundled unchanged as \texttt{sources/200616/source0.tex}.
\begin{assumption}[Max-order coverage]\label{ass:coverage}
There exists $\eta\in(0,1]$ such that along the deployed interaction (under $P^\star$ and $\pi_\theta$),
\[
\mathbb{P}^\star_{\pi_\theta}(X_t=B \mid H_{t-1}) \ge \eta\qquad\text{a.s. for all }t.
\]
\end{assumption}

\begin{assumption}[Self-contraction along AR pseudo-generation]\label{ass:self-contraction}
There exists $c_{\mathrm{sc}}\in(0,\infty)$ such that for every history value $h$ and every $s\ge 2$,
\[
\delta_s(h) \le \frac{c_{\mathrm{sc}}}{s}\, d(h),
\]
where $d(h)$ is the one-step divergence defined in \eqref{eq:def-dh} below, and $\delta_s(h)$ is the $s$-th AR conditional divergence
defined in \eqref{eq:def-delta-sh} below.
\end{assumption}

\begin{lemma}[Observed KL mismatch equals excess observed log-loss]\label{lem:delta-obs-excess-Lobs}
We have the identity
\[
\Delta^{\mathrm{obs}}_T(\theta;\pi)
=
\mathcal{L}^{\mathrm{obs}}_T(\theta;\pi)
-
\mathcal{L}^{\mathrm{obs}}_T(\star;\pi)
\;\;\ge\; 0.
\]
\end{lemma}

\begin{lemma}[Coverage turns observed mismatch into completion mismatch]\label{lem:coverage-obs-to-latent}
Assume for all $t$, $D_t\in[0,B]$ almost surely under the true environment. Moreover, also assume Assumption ~\ref{ass:coverage}. Let $\mu(\cdot\mid h)$ be any action distribution on $[0,B]$ satisfying $\mu(\{B\}\mid h)\ge\eta$.
Then
\begin{equation}\label{eq:latent-by-obs}
\mathrm{KL}(q_t^\star(\cdot\mid h)\|q_{t,\theta}(\cdot\mid h))
\;\le\;
\frac{1}{\eta}\;
\mathbb{E}_{X\sim \mu(\cdot\mid h)}
\Big[\mathrm{KL}\big(P^{\mathrm{obs}}_{h,X}\,\|\,Q^{\mathrm{obs}}_{h,X}\big)\Big].
\end{equation}
\end{lemma}

\begin{equation}\label{eq:def-dh}
d(h) := \mathrm{KL}\big(q^\star_{t,1}(\cdot\mid h)\,\|\,q_{t,1,\theta}(\cdot\mid h)\big),
\end{equation}

\begin{equation}\label{eq:def-delta-sh}
\delta_s(h)
:= \mathbb{E}_{\tilde D_{1:s-1}\sim Q^\star_t(\cdot\mid h)}
\Big[
\mathrm{KL}\big(q^\star_{t,s}(\cdot\mid h,\tilde D_{1:s-1})\,\|\,q_{t,s,\theta}(\cdot\mid h,\tilde D_{1:s-1})\big)
\Big].
\end{equation}

\begin{lemma}[Self-contraction implies completion KL is controlled by one-step KL]\label{lem:self-contraction-bound}
Assume Assumption~\ref{ass:self-contraction}. Then for every $(t,h)$,
\[
\mathrm{KL}\big(Q^\star_t(\cdot\mid h)\,\|\,Q_{t,\theta}(\cdot\mid h)\big)
\;\le\;
\bigl(1 + c_{\mathrm{sc}}\log T\bigr)\, d(h).
\]
\end{lemma}

\begin{theorem}[Relation between deployment penalty and censoring-aware predictive mismatch]
\label{thm:completion_to_ovserved}
Consider the deployed policy $\pi_\theta$ in the true environment $\P^\star$.
Assume that, for $\P^\star_{\pi_\theta}$-a.e.\ history $h$ and all $t$ the following hold:
\begin{enumerate}
\item Max-order coverage: $\P^\star_{\pi_\theta}(X_t=B\mid H_{t-1}=h)\ge \eta$, for some $\eta>0$; and
\item Self-contraction along AR trajectory: $\delta_s(h)\le (c_{\mathrm{sc}}/s)\,d(h)$, for all $s\ge 2$.
\end{enumerate}


Then the on-policy completion mismatch is upper bounded by the observed censoring-aware predictive mismatch:



\begin{equation}
\label{eq:completion_to_ovserved}
\Delta_T^{\mathrm{comp}}(\theta;\pi_\theta)
\;\le\; \frac{1 + c_{sc} \log T}{\eta}\,\Delta_T^{\mathrm{obs}}(\theta;\pi_\theta).
\end{equation}
\end{theorem}

\begin{proof}[Proof of Theorem \ref{thm:completion_to_ovserved}]
Fix $t$ and condition on $H_{t-1}=h$. The first step is to apply Lemma~\ref{lem:self-contraction-bound} and obtain
\begin{align}
\label{eq:thm3_last_1}
\mathrm{KL}\big(Q^\star_t(\cdot\mid h)\,\|\,Q_{t,\theta}(\cdot\mid h)\big)
\le (1+c_{\mathrm{sc}}\log T)\, d(h).  
\end{align}


The second step is to recognize that $d(h)=\mathrm{KL}(q_t^\star(\cdot\mid h)\|q_{t,\theta}(\cdot\mid h))$, due to the following reasoning. The first-step AR marginal $q^\star_{t,1}(\cdot\mid h)$ is (by construction of the pseudo-generation order)
the conditional law of the current-period latent demand $D_t$ given $h$, i.e. it equals $q_t^\star(\cdot\mid h)$, and similarly
$q_{t,1,\theta}(\cdot\mid h)=q_{t,\theta}(\cdot\mid h)$. 


In the third step, letting $\mu(\cdot\mid h)$ be the deployed action distribution at $h$ (i.e. $X_t\mid H_{t-1}=h\sim \mu(\cdot\mid h)$), yields the following using Lemma~\ref{lem:coverage-obs-to-latent}
\begin{align}
\label{eq:thm3_last_2}
d(h)=\mathrm{KL}(q_t^\star(\cdot\mid h)\|q_{t,\theta}(\cdot\mid h))
\le \frac{1}{\eta}\;
\mathbb{E}_{X_t\sim \mu(\cdot\mid h)}
\Big[\mathrm{KL}\big(P^{\mathrm{obs}}_{h,X_t}\,\|\,Q^{\mathrm{obs}}_{h,X_t}\big)\Big].
\end{align}

Combining the \eqref{eq:thm3_last_1} and \eqref{eq:thm3_last_2} above we obtain,
\[
\mathrm{KL}\big(Q^\star_t(\cdot\mid h)\,\|\,Q_{t,\theta}(\cdot\mid h)\big)
\le
\frac{1+c_{\mathrm{sc}}\log T}{\eta}\;
\mathbb{E}_{X_t\sim \mu(\cdot\mid h)}
\Big[\mathrm{KL}\big(P^{\mathrm{obs}}_{h,X_t}\,\|\,Q^{\mathrm{obs}}_{h,X_t}\big)\Big].
\]

Finally, taking expectation over $H_{t-1}$ under $\P^\star_{\pi_\theta}$ and summing over $t=1,\dots,T$
\[
\Delta^{\mathrm{comp}}_T(\theta;\pi_\theta)
\le
\frac{1+c_{\mathrm{sc}}\log T}{\eta}\;
\Delta^{\mathrm{obs}}_T(\theta;\pi_\theta),
\]
as claimed. Moreover, Lemma \ref{lem:delta-obs-excess-Lobs} connects the definition of $\Delta^{\mathrm{obs}}_T(\theta;\pi_\theta)$ in the manuscript (based on excess observed log-loss) to the definition using observed KL mismatch, completing our claim in the Theorem.
\end{proof}

\end{document}
